% !TEX root = ../../main-es.tex
% chapters/es/07-discusion.tex
\chapter{Discusión, amenazas a la validez y trabajo futuro}
\criterio{oral-dominio}{20}\criterio{oral-debate}{20}
\label{cap:discusion}

\porque{Sustentación oral (EICT 40\%).}{%
  El capítulo que más pesa en la defensa. Fallo típico: inflar los claims o esconder las
  limitaciones. Hacer lo contrario: declarar lo que el diseño NO garantiza y por qué es
  aceptable.}

\section{Amenazas a la validez}
\porque{Validez interna/externa.}{Confounder más plausible: el verificador obliga a escribir mejores devices (efecto Hawthorne). La ablación (B) aísla la variable. Validez externa acotada al fragmento formalizado.}
\porque{Evidencia cuantitativa, no sólo argumental.}{%
  El GLMM de §\ref{sec:protocolo} da una medida directa del confounder: la varianza del
  efecto aleatorio por tarea, $\sigma_t^2$, frente a la magnitud de los coeficientes de
  brazo. Si $\sigma_t^2$ domina, la diferencia entre brazos es pequeña comparada con la
  heterogeneidad entre tareas y el claim debe atenuarse. Si el contraste por permutación
  y el GLMM divergen, la discrepancia se reporta aquí, no se esconde.}

\section{Limitaciones}
\porque{Honestidad.}{(i) Hot-run fuera de alcance; (ii) cross-book con targets computados fuera; (iii) auto-modificación fuera (M3); (iv) medición restringida a M1 (degradación por profundidad).}

\section{Trabajo futuro}
\porque{Proyección.}{M2 (incompatibilidad de contexto), M3 (seguridad de auto-modificación — la dirección original de \OASys{}), hot-run, mecanización Lean del teorema \ref{thm:noi}.}
